
\documentclass[final,5p,times]{elsarticle}

\usepackage[utf8]{inputenc}
\usepackage[T1]{fontenc}
\usepackage{amsmath,amssymb,amsfonts}
\usepackage{booktabs}
\usepackage{graphicx}
\usepackage{xcolor}
\usepackage{url}
\usepackage{hyperref}
\IfFileExists{stfloats.sty}{\usepackage{stfloats}}{}
\usepackage{array}
\usepackage{tabularx}
\usepackage{longtable}
\usepackage{multirow}
\usepackage{enumitem}
\usepackage{makecell}
\usepackage{algorithm}
\usepackage{algorithmic}
\usepackage{balance}

\journal{Sustainable Energy, Grids and Networks}
\graphicspath{{figures/}{./}}
\biboptions{sort&compress}

\newcommand{\DeltaT}{\Delta t}
\newcommand{\EV}{\mathrm{EV}}
\newcommand{\BESS}{\mathrm{BESS}}
\newcommand{\DA}{\mathrm{DA}}
\newcommand{\RT}{\mathrm{RT}}
\newcommand{\GI}{\mathrm{GI}}
\newcommand{\PV}{\mathrm{PV}}
\newcommand{\load}{\mathrm{load}}
\newcommand{\TOU}{\mathrm{TOU}}
\newcommand{\PD}{\mathrm{PD}}
\newcommand{\NCD}{\mathrm{NCD}}
\newcommand{\SOC}{\mathrm{SOC}}
\newcommand{\soc}{\mathrm{soc}}
\newcommand{\RU}{\mathrm{RU}}
\newcommand{\RD}{\mathrm{RD}}
\newcommand{\SP}{\mathrm{SP}}
\newcommand{\NSP}{\mathrm{NSP}}
\newcommand{\WM}{\mathrm{WM}}
\newcommand{\LMP}{\mathrm{LMP}}
\newcommand{\NWM}{\mathrm{NWM}}
\newcommand{\ch}{\mathrm{ch}}
\newcommand{\dch}{\mathrm{dch}}
\newcommand{\up}{\mathrm{up}}
\newcommand{\down}{\mathrm{down}}
\newcommand{\tot}{\mathrm{tot}}
\newcommand{\base}{\mathrm{base}}
\newcommand{\forecast}{\mathrm{forecast}}
\newcommand{\real}{\mathrm{real}}
\newcommand{\executed}{\mathrm{executed}}
\newcommand{\req}{\mathrm{req}}
\newcommand{\net}{\mathrm{net}}
\newcommand{\energy}{\mathrm{energy}}
\newcommand{\capacity}{\mathrm{capacity}}
\newcommand{\avail}{\mathrm{avail}}
\newcommand{\impl}{\mathrm{impl}}
\newcommand{\deploy}{\mathrm{deploy}}
\newcommand{\pen}{\mathrm{pen}}
\newcommand{\calH}{\mathcal{H}}
\newcommand{\calT}{\mathcal{T}}
\newcommand{\calV}{\mathcal{V}}
\newcommand{\calK}{\mathcal{K}}
\newcommand{\calS}{\mathcal{S}}
\newcommand{\calD}{\mathcal{D}}
\newcommand{\calI}{\mathcal{I}}
\newcommand{\pos}[1]{\left[#1\right]_+}
\providecommand{\sep}{; }

% Safe figure command: the manuscript still compiles while figures are being regenerated.
\newcommand{\MaybeFigure}[2]{%
\IfFileExists{#1}{\includegraphics[#2]{#1}}{%
\fbox{\parbox[c][0.19\textheight][c]{0.86\linewidth}{\centering\small Missing figure file.\\ Replace with the paper-quality figure generated from the plotting notebook.}}}}

\begin{document}

\begin{frontmatter}

\title{Model Predictive Value Stacking for Campus EV Charging, Battery Storage, and Retail--Wholesale Market Participation}

\author[ucsd]{Yizhan Gu}
\author[ucsd]{Jan Kleissl\corref{cor1}}
\ead{jkleissl@ucsd.edu}
\author[totalenergies]{Hadi Eshraghi}
\author[totalenergies]{Wente Zeng}
\cortext[cor1]{Corresponding author.}
\fntext[emailnote]{Author e-mails: Yizhan Gu (yizhangu@ucsd.edu), Jan Kleissl (jkleissl@ucsd.edu), Hadi Eshraghi (hadi.eshraghi@external.totalenergies.com), and Wente Zeng (wente.zeng@totalenergies.com).}

\affiliation[ucsd]{organization={Center for Energy Research, University of California San Diego},
            city={La Jolla},
            postcode={92093},
            state={CA},
            country={USA}}

\affiliation[totalenergies]{organization={TotalEnergies},
            city={Houston},
            state={TX},
            country={USA}}

\begin{highlights}
\item EV and BESS value stacking is modeled under retail--wholesale coupling.
\item Shrinking and rolling MPC are compared using common constraints and accounting.
\item Perfect and persistence EV forecasts separate information value from control value.
\item The objective co-models EV payments, retail charges, energy, and AS revenues.
\item Unidirectional EV charging is treated as schedulable demand flexibility.
\end{highlights}

\begin{abstract}
Large workplace electric vehicle (EV) charging sites are emerging behind-the-meter distributed energy resources. Even without vehicle-to-grid export, a plugged-in EV fleet can provide operational flexibility by shifting charging in time while still satisfying departure energy requirements. This paper develops a detailed value-stacking framework for campus EV charging coordinated with battery energy storage (BESS), retail time-of-use and demand charges, day-ahead and real-time wholesale energy prices, ancillary-service capacity revenues, and EV user payments. The study compares two execution frameworks---shrinking-horizon MPC and rolling-horizon MPC---and two EV-information settings---perfect EV information and persistence-based EV forecasting. Wholesale prices and retail tariffs are treated as known exogenous inputs in the present simulations, so the forecast comparison focuses on EV availability and charging-service uncertainty. For arrived EVs, realized session information is revealed and used by the controller; for future EVs, the controller uses either perfect information or persistence-derived estimates. A July 2025 UCSD workplace charging case study with 15-minute resolution is used to evaluate the retail--wholesale tradeoff. The current setting contains daytime workplace sessions with essentially no overnight charging, so rolling MPC is not expected to systematically outperform shrinking MPC. Instead, the rolling framework is included because it is more adaptable to future scenarios with overnight sessions, cross-day flexibility, and stronger forecast uncertainty. Results are interpreted using execution-aware accounting, where gross wholesale revenue is evaluated together with retail charges, BESS feasibility, and completed EV service.
\end{abstract}

\begin{keyword}
Electric vehicle charging \sep Battery energy storage \sep Model predictive control \sep Value stacking \sep Ancillary services \sep Demand charges \sep Behind-the-meter resources
\end{keyword}

\end{frontmatter}

\section{Introduction}
\label{sec:introduction}

Workplace EV charging is often treated as a new source of peak load, but at sufficient scale it is also a flexible behind-the-meter resource. A workplace charging operator knows when vehicles are plugged in, observes how much energy has already been delivered, and can decide how fast each connected vehicle should be charged. The site normally cannot export energy from EV batteries when vehicle-to-grid (V2G) hardware is not available. Nevertheless, EV charging is not a passive load. If a vehicle will remain connected for several hours, its charging can be delayed, advanced, or reduced relative to a baseline without violating the final user energy requirement. At the fleet level, this means that unidirectional EV charging can provide a time-varying flexible demand envelope.

This flexibility becomes economically interesting when the EV fleet is co-located with a BESS. The BESS can charge or discharge directly, while the EV fleet can change the timing of its demand. Together they form a behind-the-meter portfolio whose schedule affects both a customer retail bill and potential wholesale-market revenues. The retail side includes time-of-use (TOU) energy charges, monthly peak-demand charges, and non-coincident demand charges. The wholesale side includes day-ahead (DA) and real-time (RT) energy schedules and ancillary-service (AS) products. In the California Independent System Operator (CAISO) context, ancillary services such as regulation up, regulation down, spinning reserve, and non-spinning reserve are capacity products that compensate resources for keeping deliverable operating capability available. A site can therefore earn revenue by offering capacity, but that capacity is only valuable if it is physically feasible under EV availability, EV energy requirements, BESS state of charge (SOC), and power limits.

The coupling between retail and wholesale value streams is the central difficulty. A dispatch that looks profitable from the wholesale side can increase grid import during high retail-price intervals, create a new monthly demand-charge peak, or consume BESS SOC that would otherwise reduce the retail bill. Conversely, a schedule that aggressively minimizes the retail bill may leave little headroom for wholesale participation. It is therefore not sufficient to report gross wholesale revenue. The proper objective is net behind-the-meter value, including EV user payments, retail tariff charges, wholesale energy revenue, ancillary-service revenue, and the feasibility costs associated with executing a market schedule in real time.

Forecasting adds another layer. Workplace EV sessions are uncertain because arrival times, departure times, session energy demand, and the number of vehicles vary day by day. A perfect-forecast study can be useful as an information benchmark, but it is not directly implementable. A persistence forecast, which reuses recent historical patterns, is simple and realistic, but can be wrong when daily charging behavior changes. The control algorithm must decide which forecasted sessions are available, which decisions are committed, and which decisions are re-optimized as time advances. These questions are especially important when comparing shrinking-horizon and rolling-horizon MPC. Shrinking MPC naturally closes the daily optimization horizon as the day progresses. Rolling MPC keeps a moving look-ahead window and is more flexible for scenarios with cross-day flexibility, but in a pure daytime workplace charging case with zero overnight sessions it is not expected to perform better. This point is important: the goal of the paper is not to show that rolling MPC dominates, but to show how both controller structures behave under the same value-stacking model.

This paper studies a campus-scale EV-BESS system at the University of California San Diego (UCSD). The current implemented results use EV charging and BESS dispatch. PV and building-load data exist in the broader project data pipeline, and the system diagram includes them as future model inputs, but they are not yet included in the reported optimization results. The case study focuses on July 2025, with 15-minute intervals, EV user payment revenue, retail tariff charges, DA/RT energy prices, AS capacity prices, and BESS constraints. The analysis is intentionally execution-aware: financial results are evaluated after the controller implements its decisions, not merely as a perfect-foresight planning objective.

The paper contributes the following:
\begin{enumerate}[leftmargin=*]
\item It formulates a behind-the-meter EV-BESS value-stacking model that jointly accounts for EV user payments, TOU energy charges, peak and non-coincident demand charges, DA/RT energy positions, and multiple AS capacity products.
\item It compares shrinking-horizon and rolling-horizon MPC under a common physical and economic model, avoiding the misleading interpretation that rolling MPC must be superior in every EV charging setting.
\item It evaluates both perfect and persistence EV forecasts. Perfect forecasts provide a benchmark for information value, while persistence forecasts represent a practical low-complexity forecast.
\item It explains the EV preprocessing, baseline construction, market-price processing, DA/RT execution, and representative-day operational behavior needed for reproducible campus-scale value-stacking studies.
\item It identifies a practical retail--wholesale tradeoff: higher wholesale revenue can be offset by higher behind-the-meter retail exposure, especially when EV sessions are daytime workplace sessions with no overnight flexibility.
\end{enumerate}

\section{Related Work}
\label{sec:related_work}

\subsection{EV charging flexibility without V2G}
Coordinated EV charging has been studied for distribution-grid impact reduction, valley filling, renewable integration, and charging-cost minimization \cite{clement2010impact,sortomme2011optimal,gan2013optimal,richardson2013electric,mwasilu2014electric}. Many early studies emphasize either residential charging or grid impacts from uncontrolled charging. In contrast, workplace charging has a distinct temporal structure: vehicles often arrive in the morning, remain connected for several hours, and depart before evening or late afternoon. This creates intraday flexibility but may provide little or no overnight storage-like behavior. The flexibility is therefore not equivalent to a stationary battery.

The absence of V2G does not make EVs irrelevant to grid services. A unidirectional EV charger cannot inject power into the grid, but it can reduce its charging relative to an expected baseline. From the grid perspective, reducing a controllable load can act similarly to upward flexibility, because the net load seen by the meter decreases relative to what it would otherwise have been. Increasing charging when vehicles are connected can provide downward flexibility, because the site can absorb more power. This load-modulation interpretation connects EV charging to demand response and controllable-load theory \cite{callaway2011achieving,siano2014demand}. However, any EV flexibility model must enforce driver service: each vehicle must receive its requested energy by departure, and the charging power of each session is bounded by the charger rating and the plug-in time.

\subsection{BESS value stacking and behind-the-meter tariffs}
BESS operation is frequently optimized for arbitrage, peak shaving, renewable integration, and reserve provision \cite{parisio2014model,wang2018predictive,ghosh2022battery,liu2022battery}. Behind the meter, the economic signals differ from those faced by a front-of-the-meter merchant battery. A customer-side BESS interacts with a retail tariff, so its operation changes TOU energy charges and demand charges. Demand charges are particularly important because they depend on the maximum grid import over a billing period rather than on a simple sum over time. A BESS that earns wholesale revenue during one interval can inadvertently increase a monthly peak in another interval if it charges at the wrong time.

Value stacking refers to using the same asset portfolio to earn or save money from multiple value streams. This is attractive but difficult because value streams are not independent. In the EV-BESS setting, the BESS may be asked to reduce peak demand, support EV service, arbitrage energy prices, and provide AS capacity. The EV fleet may be asked to charge cheaply, meet driver demand, absorb energy when downward flexibility is valuable, and reduce load when upward flexibility is valuable. A mathematical model must prevent simultaneous infeasible use of the same power and energy headroom. For this reason, the paper keeps the BESS physical SOC constraints, EV service constraints, and AS capacity envelopes in one optimization problem rather than layering a market schedule on top of a retail optimizer.

\subsection{Wholesale energy and ancillary services}
Wholesale energy-market participation can generate revenue by buying or selling energy positions across DA and RT prices. AS participation can generate capacity revenue by reserving upward or downward capability. Regulation up and spinning/non-spinning reserve correspond to upward capability; regulation down corresponds to downward capability. In a physical asset, upward capability may require discharge headroom or load-reduction capability, while downward capability may require charge headroom or load-increase capability. For a behind-the-meter EV-BESS system, both the BESS and the EV fleet contribute to these envelopes, but in different ways. The BESS can charge or discharge, while EVs can only increase or decrease charging relative to their baseline or charging power.

Prior studies have considered reserve provision from flexible loads and EV aggregations \cite{thrane2024reserve}. The present study differs in two practical respects. First, the EV fleet is embedded behind a retail meter with demand charges; this means a wholesale-optimal schedule may be economically unattractive after retail accounting. Second, the controller is evaluated under execution constraints. DA commitments, RT updates, and the realized EV arrivals are not merely theoretical; they determine what can actually be implemented.

\subsection{Forecasting and MPC execution}
MPC is a natural method for EV-BESS scheduling because it repeatedly solves a finite-horizon optimization problem, implements the current decision, and updates the state and forecasts. Previous microgrid and EV management studies have shown that MPC can coordinate DERs under uncertainty \cite{parisio2014model,wang2018predictive,li2021learning}. However, an MPC paper should carefully define what information is available at each time, what is forecasted, and what is committed. Otherwise, a planning result may use information that would not be available in real time.

This study explicitly separates perfect and persistence forecasts. Perfect forecasts reveal how the controller behaves when future EV sessions are known. Persistence forecasts approximate implementable operation by using recent historical sessions or recent weekday patterns. The distinction matters because EV flexibility is created by the difference between required future energy and current available charging time. Forecasting too many vehicles may overstate flexibility, while forecasting too few vehicles may underutilize flexibility and reduce market value.

\subsection{Prior UCSD and collaborator work}
Two predecessor studies by Chen and collaborators provide important context for the present work \cite{chen2024costoptimal,chen2026realtime}. They are cited here as prior UCSD/collaborator work on EV charging data processing, charging-control motivation, and campus-scale flexibility analysis. They should not be read as previous BESS or wholesale-market formulations. In particular, the present paper does not claim that those studies already modeled BESS dispatch, DA/RT wholesale participation, or ancillary-service products. Rather, they motivate the campus EV charging setting and the need to move from descriptive charging analysis toward operational control.

The contribution of the present paper is to extend that research line in three directions. First, EV charging is coordinated with a BESS, so the EV fleet is no longer treated as an isolated flexible load. Second, the objective combines EV user payments, retail energy and demand charges, wholesale energy revenue, and AS capacity revenue in one accounting framework. Third, the paper compares shrinking and rolling MPC under perfect and persistence EV-information settings. This creates a more explicit bridge between EV session preprocessing, forecast construction, MPC execution, and financial evaluation.

The predecessor papers are therefore used as background and motivation, while the equations and results in this manuscript follow the current simulation implementation and the cleaned formulation. Where details differ, the current EV/BESS/market implementation and the present mathematical model are treated as the authoritative description.
\subsection{Research gap}
The literature motivates the components of the problem---EV flexibility, BESS operation, MPC, and market participation---but the combination remains challenging. A practical campus operator needs to know not only whether wholesale participation creates gross revenue, but also whether it improves the net value after retail charges and EV user payments. The operator also needs to know whether a more adaptive rolling controller is actually better for a daytime workplace data set, or whether a shrinking daily controller is more appropriate when all sessions close within the day. This paper addresses that gap by formulating and evaluating an execution-aware EV-BESS value-stacking controller with both forecast and controller comparisons.

\section{System Description and Data Flow}
\label{sec:system_description}

\subsection{Behind-the-meter site}
The modeled site is a campus EV charging system with a co-located BESS behind a retail meter. The EV chargers serve workplace sessions. The BESS is controlled by the site operator and can be scheduled for retail peak management, energy arbitrage, and AS capacity. The retail meter observes the net grid import after EV charging and BESS dispatch. In the broader project architecture, building load and PV data are also available. The present reported simulations do not yet include those streams as active optimized terms, so the current results should be interpreted as an EV-BESS focused stage of a larger DER platform.

The control problem has three layers. The physical layer includes connected EVs, BESS SOC, and power limits. The retail layer maps net meter import to TOU and demand charges. The wholesale layer maps energy and AS decisions to revenue. A correct value-stacking model must couple all three layers. For example, offering regulation down may require the site to absorb additional energy, which may be feasible only if there is BESS charge headroom or connected EV charging headroom. That same absorption can increase retail import and possibly affect demand charges.

\begin{figure*}[!t]
\centering
\MaybeFigure{system_architecture.png}{width=0.88\textwidth}
\caption{Behind-the-meter EV--BESS value-stacking architecture. The active July 2025 simulations coordinate workplace EV charging and BESS dispatch through an MPC controller that sees EV information, tariff parameters, wholesale prices, AS prices, BESS state, and DA/RT commitments. PV and building-load data streams are shown as available extensions but are not included in the reported EV--BESS optimization results.}
\label{fig:system_architecture}
\end{figure*}

\subsection{Data sources and preprocessing logic}
The case-study pipeline uses several data streams. EV session data are converted into 15-minute interval representations. Each session is associated with an arrival time, departure time, delivered or requested energy, and charging power limit. The preprocessing removes or flags sessions that cannot be represented consistently at the 15-minute control resolution, converts timestamps into the simulation time zone, and constructs per-interval availability matrices. The implementation uses a maximum charger power of 6.6 kW per active EV session in the aggregate capability calculation. The EV user-payment revenue is computed from delivered EV energy and a user charging price; the current results imply approximately 54.87 MWh of delivered July EV energy when using a 0.40 \/kWh charging price.

The baseline pipeline creates the EV baseline $B_t^{\EV}$ used for retail and wholesale flexibility accounting. The baseline is important because the EV fleet does not export. Its market-relevant contribution is the ability to increase or decrease charging relative to a reference schedule while still meeting service. A baseline that is too high exaggerates upward flexibility; a baseline that is too low exaggerates downward flexibility. The manuscript therefore treats baseline construction as a critical modeling choice and identifies baseline revision as an ongoing validation item.

Wholesale data include DA and RT LMPs and AS prices for regulation up, regulation down, spinning reserve, and non-spinning reserve. The price-processing notebooks download or parse market files, align the data to the 15-minute simulation index, and build the price arrays used by the optimizer. AS prices enter as capacity prices, while deployment factors approximate the expected energy movement associated with awards. Retail data include TOU energy prices, peak-demand-charge rates, and non-coincident demand-charge rates. The simulation setup uses July/summer tariff parameters; seasonal refinements should be handled when the study is extended to additional months.

\subsection{Workflow overview}
Figure~\ref{fig:workflow} summarizes the publication-quality control workflow. The schematic links raw data, EV-information construction, MPC optimization, first-step execution, state update, result logging, and financial accounting. The manuscript also includes tabular and algorithmic descriptions so the control logic is clear even when the workflow diagram is viewed at reduced journal scale.

\begin{figure*}[!t]
\centering
\MaybeFigure{market_mpc_workflow.png}{width=0.90\textwidth}
\caption{Forecasting and MPC workflow. The diagram distinguishes EV-information construction, DA schedule generation, RT execution, shrinking-horizon updates, rolling-horizon updates, and ex-post financial accounting.}
\label{fig:workflow}
\end{figure*}

\section{Forecasting Methods}
The controller requires forecasts of EV charging demand because future EV arrivals, departure times, and requested energy determine the feasible EV charging envelope. In the present simulations, the uncertain quantities are associated with EV sessions. Retail tariff schedules and wholesale-market prices are treated as known exogenous inputs over the optimization horizon. This modeling choice isolates the value of EV information and EV forecast quality from price-forecasting questions. It also matches the current simulation implementation, where CAISO day-ahead and real-time price series and AS price series are read from processed market data files rather than predicted by a price model.

\subsection{Forecast objects}
At each control time $t$, the EV forecast object describes a set of currently connected and future EV sessions over the prediction horizon $\calH_t$. For an EV $i$, the relevant information includes its arrival interval $a_i$, departure interval $d_i$, requested energy $E_i^{\req}$, maximum charging power $\bar p_i^{\EV}$, and availability indicator $b_{t,i}^{\EV}$. Once an EV has arrived, the controller observes its realized session information and uses that information in subsequent optimizations. For EVs that have not yet arrived, the controller either uses perfect information for benchmarking or persistence-based estimates for an implementable forecast.

This distinction is important. In this study, ``perfect'' does not mean perfect electricity-price prediction; the price series are not forecasted. Instead, perfect forecasting means that future EV session information is made available to the optimizer as an information benchmark. Persistence forecasting means that future EV charging demand is estimated from recent historical charging behavior and then translated into future session or aggregate charging requirements.

\subsection{Perfect EV-information benchmark}
The perfect forecast case uses the realized EV session set for the study period when constructing the EV constraints. It answers the question: how well could the same optimization and execution framework perform if future EV arrivals, departures, and energy requests were known? This case is not claimed to be implementable in practice. It is an information benchmark that separates the effect of forecast uncertainty from the effect of the controller structure. A large gap between perfect and persistence cases indicates that improved EV forecasting may have economic value. A small gap indicates that the controller is less sensitive to the forecast errors present in the tested month.

Because perfect EV information is available over the horizon, shrinking MPC can plan a full daily or remaining-day schedule with fewer forecast revisions. Rolling MPC can also use perfect information, but its advantage is not expected to appear in this daytime workplace case simply because most EV sessions begin and end within the same day and overnight carryover is rare. The perfect benchmark is therefore used to interpret both forecast value and controller value.

\subsection{Persistence EV forecast}
Persistence forecasting is the main implementable baseline. The idea is that recent or representative historical charging behavior is used to estimate future workplace charging demand. In the current EV forecasting pipeline, two quantities are central: the number of EVs expected to be present and the amount of session energy expected to be requested. These estimates are converted into EV availability and service requirements over the prediction horizon. As time advances, newly arrived EVs reveal their realized information, and the forecast object is updated.

The persistence forecast is intentionally simple. Its role is not to be the final forecasting method for all future deployment cases, but to provide a transparent benchmark that can be reproduced and compared against perfect information. For a campus workplace charging site, persistence can perform reasonably because weekday charging patterns are often recurring. However, it can still misrepresent unusual days, holiday effects, or changes in arrival intensity. Those forecast errors affect the EV flexibility envelope and, through the BTM balance, can also affect BESS dispatch, retail charges, and wholesale-market participation.

\subsection{Forecast construction as an optimization input}
For both perfect and persistence cases, the output of the forecast module is not merely a scalar load profile. It is a structured set of EV constraints. The forecast determines which EVs are assumed to be available, how much energy must be delivered before departure, and what aggregate charging capability exists at each 15-minute interval. The controller then combines that EV forecast with BESS state, retail tariff parameters, known market prices, and previous market commitments. In this way, the forecast module determines feasible flexibility, while the optimization module decides how to allocate that flexibility among EV service, retail-cost reduction, and wholesale-market products.
\section{MPC Control Frameworks, Bidding Timeline, and Execution Logic}
The same EV/BESS value-stacking model is solved under two execution frameworks: shrinking-horizon MPC and rolling-horizon MPC. Both frameworks use the same device constraints, tariff accounting, wholesale products, and financial evaluation. They differ in how the horizon is updated and how future decisions are recomputed.

\subsection{Why shrinking and rolling MPC are both studied}
Shrinking-horizon MPC is a natural baseline for the July workplace charging case. At the beginning of a daily operating problem, the controller optimizes over the remaining intervals in the day. As the day progresses, the remaining horizon shrinks. This is well matched to daytime workplace charging because most sessions arrive and depart during the same day, and the current data set contains essentially no overnight sessions. For this reason, shrinking MPC is not expected to be weaker than rolling MPC in the present scenario. In several cases it can be stronger because it avoids repeatedly expanding the planning window into periods that contain little additional useful EV flexibility.

Rolling-horizon MPC solves a moving-window problem. After each 15-minute execution step, the controller advances the horizon and rebuilds the forecast object. Rolling MPC is therefore more general. It is expected to become more valuable when overnight charging sessions, residential charging, fleet depots, or multi-day uncertainty are included. Those settings create cross-day coupling and make repeated forecast updates more valuable. The present paper therefore does not present rolling MPC as a guaranteed improvement for the July workplace case. Instead, it compares shrinking and rolling as two execution structures under a common model, and it interprets the results in the context of the available EV flexibility.

\subsection{Market bidding and MPC timeline}
The physical controller must also respect the market-information timeline. Fig.~\ref{fig:bidding_timeline} summarizes the intended sequence. Historical EV/PV/building data are collected before the operating day. In the current reported model, the EV stream is active while PV and building-load streams are retained for later integration. A day-ahead forecast is created on Day--1 and used to solve the DA optimization. The resulting DA wholesale bid is submitted before the DA market deadline. After the DA market clears, awarded DA quantities become commitments or references for the RT controller. During Day 0, the RT controller updates every 15 minutes using realized EV arrivals, updated BESS SOC, known price streams, and the remaining DA commitments. It then solves an RT problem and executes only the current 15-minute dispatch.

\begin{figure*}[!t]
\centering
\MaybeFigure{market_bidding_mpc_timeline.png}{width=0.94\textwidth}
\caption{Market bidding and MPC execution timeline. The DA stage uses Day--1 information to create a forecast, solve the DA optimization, and submit DA energy/AS bids. After DA market clearing, awarded DA quantities are passed to the RT controller. On Day 0, EV data and forecasts are updated every 15 minutes, the RT MPC solves the current dispatch problem, and only the first 15-minute decision is executed.}
\label{fig:bidding_timeline}
\end{figure*}

The DA/RT distinction matters because the RT problem should not freely ignore DA quantities that have already been submitted or awarded. The current simulation implementation therefore uses a penalty-based RT reference-tracking formulation rather than a hard equality formulation. A hard formulation can become infeasible in rolling simulations if future DA commitments are inconsistent with newly observed EV availability, BESS SOC, or remaining flexibility. The penalty formulation keeps the RT problem feasible while discouraging unnecessary deviations from the DA reference. This is especially important for month-long simulations where a single infeasible RT step would stop the scenario.

\subsection{Day-ahead and real-time roles}
The DA optimization determines the schedule and submitted wholesale position for the next operating day. It uses the selected EV forecast, known tariff parameters, known DA price series, and AS price series. The RT optimization is executed at 15-minute resolution. It observes realized states, updates the EV forecast, tracks relevant DA values with a penalty, and computes the immediate EV/BESS dispatch and RT market quantities. Financial results are computed from executed decisions rather than from future decisions that may later be overwritten.

\subsection{Example: execution at 9:00 AM and 9:15 AM}
Consider Day 0 at 9:00 AM. At this time, the controller knows the EVs that have already arrived and observes their realized session information. It also knows the current BESS SOC and accumulated retail demand-charge states. The EV forecast module fills in future EV sessions over the remaining horizon using either perfect EV information or persistence estimates. The price inputs are read from the processed market and tariff data, not forecasted by the EV module. The MPC then solves the MILP with EV service constraints, BESS SOC constraints, BTM meter balance, retail-cost terms, wholesale revenue terms, AS capability constraints, and DA reference-tracking penalties.

Only the 9:00--9:15 AM EV and BESS decisions are implemented. The rest of the optimized trajectory is a plan, not an executed schedule. At 9:15 AM, delivered EV energy and BESS SOC are updated, any newly arrived EVs reveal their session information, and the forecast object is rebuilt. Shrinking MPC solves over the remaining part of the day, whereas rolling MPC shifts its moving horizon forward. The cycle repeats until the operating horizon ends.

\begin{table*}[!t]
\centering
\caption{Execution logic for the 9:00 AM to 9:15 AM update. The same physical model is used by both MPC frameworks; only the horizon update and forecast replacement rule differ.}
\label{tab:execution_example}
\scriptsize
\begin{tabularx}{\textwidth}{@{}p{0.16\textwidth}X X@{}}
\toprule
Step & Shrinking-horizon MPC & Rolling-horizon MPC \\
\midrule
At 9:00 AM & Build a forecast and optimize over the remaining intervals of Day 0. & Build a forecast and optimize over the rolling window that begins at 9:00 AM.\\
EV information & Arrived EVs use realized session information; future EVs use perfect or persistence EV information. & Same EV information rule, but the forecast window shifts forward at the next step.\\
Market information & DA commitments and known price streams enter the RT problem; deviations are penalized. & Same DA/RT tracking logic.\\
Execution & Implement only the 9:00--9:15 AM decisions. & Implement only the 9:00--9:15 AM decisions.\\
At 9:15 AM & Remove the executed interval and optimize the remaining day. & Shift the window to start at 9:15 AM and rebuild the forecast object.\\
Expected role & Strong baseline for daytime workplace charging without overnight carryover. & More adaptable for future overnight, multi-day, or higher-uncertainty EV settings.\\
\bottomrule
\end{tabularx}
\end{table*}

\subsection{Scenario definitions}
The experimental design is a two-by-two controller and EV-information comparison, crossed with retail-only and wholesale-enabled operation. The controller dimension compares shrinking-horizon MPC with rolling-horizon MPC; the EV-information dimension compares perfect EV information with a persistence EV forecast; and the operating-mode dimension compares retail-only scheduling with retail+wholesale value stacking. Retail-only cases disable wholesale energy and AS participation. Retail+wholesale cases allow DA/RT energy and AS products subject to the same EV/BESS physical constraints and the same retail meter accounting. Because all cases use the same EV service requirement and financial accounting, differences in net value can be attributed to controller execution, EV information, and wholesale participation rather than to inconsistent service delivery.
\section{Optimization Formulation}
This section presents the core mathematical model in the main text. Detailed linearization, binary direction variables, AS duration constraints, and revenue-decomposition equations are provided in Appendix~A. The notation is introduced with the model components rather than placed in a separate nomenclature section, so that each symbol is defined close to where it is used.

The model is solved repeatedly under either shrinking-horizon or rolling-horizon MPC. Let $\calH_t$ denote the active prediction horizon at control time $t$. For shrinking MPC, $\calH_t$ contains the remaining intervals in the operating day. For rolling MPC, $\calH_t$ is a moving look-ahead window that shifts forward after each execution step. Both controllers use the same physical constraints, retail-meter accounting, wholesale-market variables, and EV-service requirements; they differ in how the horizon is updated and how future EV information is replaced as time advances. The controller therefore solves the value-stacking model below at each decision time, executes only the current interval, updates the realized state, and then repeats the process with the appropriate shrinking or rolling horizon.

\begin{table*}[!t]
\centering
\caption{Core notation used in the formulation. Additional product-specific variables are defined where they appear in Appendix~A.}
\label{tab:notation_core}
\scriptsize
\begin{tabularx}{\textwidth}{@{}p{0.16\textwidth}p{0.78\textwidth}@{}}
\toprule
Symbol & Meaning \\
\midrule
$t\in\calT$, $i\in\calI$ & 15-minute interval and EV/session index.\\
$\DeltaT$ & Interval length in hours.\\
$p_{t,i}^{\EV}$, $p_t^{\EV}$ & Charging power of EV $i$ and aggregate EV fleet charging power.\\
$e_{t,i}$, $E_i^{\req}$ & Delivered energy state and required departure energy of EV $i$.\\
$b_{t,i}^{\EV}$, $p_t^{\EV,\avail}$ & EV plug-in availability indicator and aggregate available charging capability.\\
$p_t^{\ch,\BESS}$, $p_t^{\dch,\BESS}$, $\soc_t^{\BESS}$ & BESS charge power, discharge power, and state of charge.\\
$p_t^{\GI}$ & Behind-the-meter grid import measured at the retail meter.\\
$C_t^{\TOU}$, $C^{\PD}$, $C^{\NCD}$ & Retail energy, peak-demand, and non-coincident demand-charge rates.\\
$C_t^{\LMP,\DA}$, $C_t^{\LMP,\RT}$ & Known DA and RT wholesale energy price inputs.\\
$c_t^{k,s}$ & AS capacity award for product $k$ and direction/status $s$.\\
$R^{\EV}$, $R^{\WM}$ & EV user-payment revenue and wholesale-market revenue.\\
\bottomrule
\end{tabularx}
\end{table*}

\subsection{Objective function}
At each optimization call, the controller minimizes net cost, equivalently maximizing net value. The compact objective is
\begin{align}
\min \quad
J = C^{\TOU}+C^{\PD}+C^{\NCD}-R^{\EV}-R^{\WM}+\rho^{\WM}+\rho^{\BESS},
\label{eq:objective_main}
\end{align}
where $C^{\TOU}$ is the retail energy charge, $C^{\PD}$ is the peak-demand charge, $C^{\NCD}$ is the non-coincident demand charge, $R^{\EV}$ is EV user-payment revenue, $R^{\WM}$ is wholesale revenue, $\rho^{\WM}$ is a DA/RT schedule-tracking penalty used in RT, and $\rho^{\BESS}$ is a small implementation penalty used to discourage unnecessary BESS cycling or ramping. This objective makes the main economic point of the paper visible: wholesale revenue appears with a negative sign in a cost-minimization objective, but it is not free. The same dispatch also changes retail import and therefore retail cost.

\subsection{Expanded objective components}
For reproducibility, the compact objective in Eq.~\eqref{eq:objective_main} can be expanded into interval and billing-period components. The interval EV revenue is
\begin{align}
r_t^{\EV}=C^{\EV}p_t^{\EV}\DeltaT.
\label{eq:interval_ev_rev}
\end{align}
The interval TOU energy cost is
\begin{align}
c_t^{\TOU}=C_t^{\TOU}p_t^{\GI}\DeltaT.
\label{eq:interval_tou}
\end{align}
The wholesale energy and AS revenue can be written as
\begin{align}
r_t^{\WM}=&\DeltaT\left(C_t^{\mathrm{LMP},\DA}p_t^{\DA}+C_t^{\mathrm{LMP},\RT}p_t^{\RT}\right)\nonumber\\
&+\DeltaT\sum_{k\in\calK}\sum_{s\in\calS}C_t^{k,s}c_t^{k,s}+r_t^{\mathrm{deploy}}.
\label{eq:interval_wm_rev}
\end{align}
Demand charges are not simple interval sums. In the implemented monthly accounting, a new demand-charge cost is incurred when the current grid import exceeds the previous billing-period maximum. An equivalent epigraph representation uses variables $z^{\PD}$ and $z^{\NCD}$:
\begin{align}
z^{\PD}&\ge p_t^{\GI},\quad t\in\calT_{\PD},\label{eq:pd_epigraph}\\
z^{\NCD}&\ge p_t^{\GI},\quad t\in\calT,\label{eq:ncd_epigraph}\\
C^{\PD}&=C^{\PD,rate}z^{\PD},\qquad C^{\NCD}=C^{\NCD,rate}z^{\NCD}.\label{eq:dem_epigraph_cost}
\end{align}
This representation is useful because it keeps the demand-charge model linear while preserving the economic meaning of the monthly maximum. It also explains why daily financial plots should be interpreted with care: a large demand-charge impact can be triggered by a single interval.

\subsection{EV charging model}
Each EV session has an arrival time $a_i$, departure time $d_i$, requested energy $e_i^{\req}$, and availability indicator $b_{t,i}^{\EV}$. The charging power is constrained by availability and charger rating:
\begin{align}
0 \le p_{t,i}^{\EV} \le \hat P^{\EV} b_{t,i}^{\EV}.
\label{eq:ev_power_bound_main}
\end{align}
The delivered energy state evolves as
\begin{align}
e_{t,i}=e_{t-1,i}+p_{t,i}^{\EV}\DeltaT/\eta_i^{\EV},
\label{eq:ev_energy_main}
\end{align}
and the service requirement is
\begin{align}
e_{d_i,i}\ge e_i^{\req}.
\label{eq:ev_service_main}
\end{align}
The aggregate EV charging power is
\begin{align}
p_t^{\EV}=\sum_{i\in\calV} p_{t,i}^{\EV}.
\label{eq:ev_aggregate_main}
\end{align}

Equation~\eqref{eq:ev_aggregate_main} is more than notation. It turns many individual charging decisions into the fleet-level demand seen by the meter and the wholesale module. Because V2G is not modeled, $p_{t,i}^{\EV}$ and $p_t^{\EV}$ are nonnegative. The EV fleet still has economic flexibility because the optimizer can decide when to deliver the required energy. If the fleet would normally charge according to a baseline $B_t^{\EV}$, then reducing $p_t^{\EV}$ below that baseline releases upward flexibility; increasing $p_t^{\EV}$ above that baseline absorbs additional power and can provide downward flexibility. The flexibility is real only if the final energy requirement \eqref{eq:ev_service_main} remains satisfied.

The dynamic aggregate charging capability is
\begin{align}
P_{t,\max}^{\EV}=\hat P^{\EV}\sum_{i\in\calV}b_{t,i}^{\EV}.
\label{eq:ev_capability_main}
\end{align}
This value changes throughout the day because vehicles arrive and depart. It determines how much EV demand can be increased at time $t$. It also appears in the AS and net-power envelopes in Appendix~\ref{app:complete_milp}. A high $P_{t,\max}^{\EV}$ does not automatically mean high flexibility: if the connected vehicles are already nearly fully charged, downward flexibility may be limited by remaining energy need; if they are close to departure, upward flexibility may be limited by service completion.

\subsection{BESS model}
The BESS has charging power $p_t^{\ch,\BESS}$, discharging power $p_t^{\dch,\BESS}$, and SOC $\soc_t^{\BESS}$. The SOC dynamics are
\begin{align}
\soc_t^{\BESS}=\soc_{t-1}^{\BESS}+\frac{\DeltaT}{C^{\BESS}}
\left(p_t^{\ch,\BESS}\sqrt{\gamma}-\frac{p_t^{\dch,\BESS}}{\sqrt{\gamma}}\right),
\label{eq:bess_soc_main}
\end{align}
with bounds
\begin{align}
\soc_{\min}^{\BESS}\le \soc_t^{\BESS}\le \soc_{\max}^{\BESS}.
\label{eq:bess_soc_bound_main}
\end{align}
The daily terminal equality is
\begin{align}
\soc_{T_d}^{\BESS}=\soc_{0,d}^{\BESS},
\label{eq:bess_terminal_main}
\end{align}
where $T_d$ denotes the end of the daily accounting horizon. This prevents the controller from artificially improving a daily objective by emptying or filling the BESS at the end of the simulated day.

The BESS is the physical resource that can actually move energy across time. It can charge when wholesale or retail conditions favor absorption, discharge when peak reduction or energy revenue is valuable, and hold SOC to provide upward or downward AS capability. It therefore complements the EV fleet. EV flexibility is constrained by user service and plug-in availability; BESS flexibility is constrained by SOC, power, and throughput.

\subsection{Behind-the-meter power balance}
The retail meter import is modeled as
\begin{align}
p_t^{\GI}=p_t^{\load}+p_t^{\EV}+p_t^{\ch,\BESS}-p_t^{\dch,\BESS}-p_t^{\PV}.
\label{eq:btm_balance_main}
\end{align}
In the reported EV-BESS results, $p_t^{\load}$ and $p_t^{\PV}$ are held out of the active optimization and are treated as future integration streams. Equation~\eqref{eq:btm_balance_main} is still included because it is the general behind-the-meter balance and because it clarifies how PV and building load will enter the extended model.

This balance is the bridge between wholesale actions and retail costs. If the BESS charges to support regulation down or energy arbitrage, $p_t^{\GI}$ can increase. If that happens during an expensive TOU period or during the monthly peak, retail costs can rise. Likewise, if EV charging is shifted to avoid a demand charge, the site may lose wholesale opportunities. This is why the model optimizes net value rather than maximizing $R^{\WM}$ alone.

\subsection{Retail tariff terms}
The TOU energy charge is
\begin{align}
C^{\TOU}=\sum_{t\in\calT} C_t^{\TOU}p_t^{\GI}\DeltaT.
\label{eq:tou_main}
\end{align}
The demand-charge terms are modeled through maximum grid-import quantities:
\begin{align}
C^{\PD}=C^{\PD,rate}\max_{t\in\calT_{\PD}} p_t^{\GI},\qquad
C^{\NCD}=C^{\NCD,rate}\max_{t\in\calT}p_t^{\GI}.
\label{eq:demand_charge_main}
\end{align}
In the simulation workflow, these max functions are represented with auxiliary epigraph variables. The peak-demand window $\calT_{\PD}$ depends on the retail tariff. The non-coincident demand charge applies to the maximum import over the broader billing period. These max terms are why a single interval can have a large monthly financial impact.

\subsection{Wholesale revenue and AS products}
Wholesale revenue includes DA and RT energy revenue and AS capacity revenue:
\begin{align}
R^{\WM}=&\sum_{t\in\calT}\DeltaT\left(C_t^{\mathrm{LMP},\DA}p_t^{\DA}+C_t^{\mathrm{LMP},\RT}p_t^{\RT}\right)\nonumber\\
&+\sum_{t\in\calT}\DeltaT\sum_{k\in\calK}\sum_{s\in\calS} C_t^{k,s}c_t^{k,s} + R^{\mathrm{deploy}},
\label{eq:wm_revenue_main}
\end{align}
where $R^{\mathrm{deploy}}$ accounts for expected deployment energy using product-specific deployment factors. In the simulation workflow, regulation products are assigned larger deployment factors than spinning and non-spinning reserve. Upward products and downward products have opposite signs in the deployment-energy accounting.

The capacity variables $c_t^{k,s}$ are not unconstrained bids. They are bounded by EV and BESS power headroom, BESS SOC duration requirements, EV baseline deviations, and directionality constraints. For example, the combined upward capability includes BESS discharge headroom plus the ability to reduce EV charging relative to baseline. The combined downward capability includes BESS charge headroom plus the ability to increase EV charging up to $P_{t,\max}^{\EV}$. The complete envelopes are listed in Appendix~\ref{app:complete_milp}.

\subsection{Why EV flexibility and BESS dispatch must be co-optimized}
A tempting simplification is to first optimize EV charging for user service and then optimize the BESS for market value. That decomposition is not valid for this paper's value-stacking problem because EV charging and BESS dispatch share the same meter, the same retail demand-charge state, and the same wholesale feasibility envelope. If EV charging is delayed to reduce meter import during a peak-demand interval, the BESS may have more opportunity to charge or may need to discharge less. If the BESS charges to support a downward product, the EV fleet may need to reduce charging to avoid exceeding a demand-charge threshold. These interactions appear directly in the meter balance \eqref{eq:btm_balance_main} and the capability constraints in Appendix~\ref{app:complete_milp}.

The formulation also avoids double counting EV flexibility. EVs can only be used if the resulting charging schedule still delivers the required energy by departure. In other words, EV flexibility is a scheduling margin, not a new energy source. The model can reduce charging now only if it can charge later before departure. The aggregate baseline and capability variables make this margin visible to the wholesale module, while the individual EV energy constraints protect user service.

\subsection{Implementation constraints and RT reference tracking}
The DA and RT problems share the same physical EV/BESS model, but RT operation has one additional practical requirement: it should remain close to previously submitted DA quantities when those quantities are active. Enforcing all such DA quantities as hard equalities can make a rolling month-long simulation infeasible because the rolling horizon may later observe EV arrivals, BESS SOC, or remaining flexibility that are inconsistent with a previously value. For this reason, the present implementation uses penalty-based reference tracking:
\begin{equation}
\Pi_t^{\RT}
=
\sum_{\tau\in\calH_t}
\left(
\rho_p |p_{\tau}^{\RT}-p_{\tau}^{\DA,\mathrm{ref}}|
+
\sum_{k,s}\rho_c |c_{\tau}^{k,s}-c_{\tau}^{k,s,\mathrm{ref}}|
\right).
\end{equation}
The penalties discourage unnecessary deviations from DA references while preserving feasibility. This choice is particularly important for rolling MPC, where a strict hard-commitment formulation can become infeasible over long monthly simulations. Shrinking MPC can sometimes tolerate stricter DA commitment constraints because the daily horizon is shorter and does not repeatedly expand into future windows, but using a common penalty-based implementation improves comparability across controllers. The penalty terms should therefore be interpreted as numerical and operational regularization, not as an additional source of revenue.
\subsection{EV user payments and service economics}
EV user revenue is modeled as
\begin{align}
R^{\EV}=\sum_{t\in\calT} C^{\EV}p_t^{\EV}\DeltaT.
\label{eq:ev_revenue_main}
\end{align}
This term represents user payment for delivered charging energy. It does not mean that the optimizer should overcharge vehicles; session requirements and upper energy bounds prevent that. In the July 2025 accounting, EV user revenue is the same across completed cases when the same EV service is delivered. The differences in net value therefore come from retail costs and wholesale revenue rather than from changing total EV energy served.

\section{Case Study}
\label{sec:case_study}

\subsection{UCSD workplace charging setting}
The case study represents a UCSD campus workplace charging data set for July 2025. The simulation interval is 15 minutes. The current July 2025 result set focuses on EV charging and a 250 kW / 332 kWh BESS with SOC bounds of 0.05--0.95 and round-trip efficiency parameter $\gamma=0.9$. The EV charger limit used in the aggregate capability calculation is 6.6 kW per connected EV session. The solver configuration uses CVXPY and Gurobi with a 0.1\% MIP gap, 10 solver threads, and a 0.5 s RT time limit. These values are implementation settings rather than universal recommendations.

\begin{table*}[!t]
\centering
\caption{Case-study parameters used in the current simulation implementation. Values correspond to the July 2025 EV--BESS simulation setting.}
\label{tab:case_parameters}
\scriptsize
\begin{tabularx}{\textwidth}{@{}p{0.23\textwidth}p{0.18\textwidth}X@{}}
\toprule
Parameter & Value & Notes \\
\midrule
Simulation month & July 2025 & Current sample-month study; additional months are future work.\\
Time resolution & 15 min & One day has 96 intervals.\\
BESS power rating & 250 kW & Current simulation setting.\\
BESS energy capacity & 332 kWh & Current simulation setting.\\
BESS SOC bounds & 0.05--0.95 & Daily terminal equality used for daily accounting.\\
BESS efficiency parameter & $\gamma=0.9$ & Implemented with $\sqrt{\gamma}$ charge/discharge split.\\
EV charger rating & 6.6 kW/session & Used for aggregate EV capability.\\
EV charging price & 0.40 \/kWh & Used for EV user-payment revenue.\\
Peak-demand rate & 3.05 \/kW & July/summer tariff setting used in the current simulation.\\
Non-coincident demand rate & 15.38 \/kW & July/summer tariff setting used in the current simulation.\\
Solver & Gurobi via CVXPY & MILP implementation.\\
MIP gap & $10^{-3}$ & Current simulation setting.\\
RT time limit & 0.5 s & Used for fast repeated RT execution.\\
PV/building load & Not active in reported results & Data available; integration is future work.\\
\bottomrule
\end{tabularx}
\end{table*}

\subsection{EV data processing}
The EV preprocessing converts raw charging-session records into the session-level and interval-level objects required by the optimizer. The session-level data include arrival time, departure time, delivered energy, and driver/session identifiers when available. The interval-level data define which vehicles are connected in each 15-minute period and how much charging power can be assigned. Sessions that cannot be represented consistently at the simulation resolution should be filtered or flagged, because infeasible session data can cause artificial unmet energy or unrealistic flexibility.

The key modeling choice is that the optimizer controls charging power, not session arrival. Once a vehicle is observed, its current availability is known. Future vehicles are forecast. In a perfect forecast, future arrivals and energy requirements are known from the realized data. In a persistence forecast, future arrivals and session energy are inferred from recent patterns. The preprocessing therefore feeds both the realized execution state and the forecast set used by the MPC.

\subsection{EV baseline construction}
The EV baseline $B_t^{\EV}$ is used to define load deviation and flexibility. In the current simulation workflow, the baseline is generated from a baseline simulation and converted from interval energy to power. This baseline enters the EV net charge/discharge decomposition in Appendix~\ref{app:complete_milp}. It is especially important for AS products, because load reduction or load increase is measured relative to an expected charging trajectory.

The baseline is not merely cosmetic. A baseline determines how much upward or downward flexibility the EV fleet is allowed to claim. If baseline charging assumes all connected EVs charge immediately at maximum power, the fleet may appear to provide large upward flexibility by reducing charging. If the baseline is conservative, downward flexibility may be smaller. Because of this sensitivity, baseline construction is identified as a future validation item. The current paper reports results using the implemented baseline and avoids claiming that this baseline is unique or final.

\subsection{Market-price processing}
The DA/RT LMP and AS price pipelines align market data to the same 15-minute simulation intervals used by EV and BESS control. DA and RT energy prices enter the energy-revenue term. AS prices enter the capacity-revenue term for regulation up, regulation down, spinning reserve, and non-spinning reserve. The implementation also uses product deployment factors to account for expected energy movement associated with awards. The price-processing notebooks therefore play two roles: they create the arrays used in the optimizer, and they create daily diagnostic files that can be used to interpret representative-day behavior.

\begin{table*}[!t]
\centering
\caption{Data-processing modules used to create the EV, price, market, dispatch, and accounting inputs.}
\label{tab:data_pipeline}
\scriptsize
\begin{tabularx}{\textwidth}{@{}p{0.20\textwidth}p{0.22\textwidth}p{0.22\textwidth}X@{}}
\toprule
Pipeline component & Main inputs & Main outputs & Manuscript relevance \\
\midrule
EV post-processing & Raw charging-session data, timestamps, session energy, charger information & Clean session table, 15-min availability, arrival/departure indices, delivered/requested energy & Defines $a_i$, $d_i$, $e_i^{\req}$, $b_{t,i}^{\EV}$, and the fleet capability envelope.\\
EV baseline generation & Processed sessions and baseline simulation settings & $B_t^{\EV}$ and interval baseline energy/power & Determines upward/downward EV flexibility relative to the baseline.\\
Persistence EV forecast & Recent historical session energy and EV counts & Forecast session-energy and EV-count profiles & Main implementable EV forecast method.\\
Perfect EV forecast & Realized future session information & Information-benchmark EV forecast & Upper-bound/reference forecast case.\\
LMP processing & CAISO DA/RT market files & Interval-aligned $C_t^{\mathrm{LMP},\DA}$ and $C_t^{\mathrm{LMP},\RT}$ & Drives DA/RT energy revenue.\\
AS processing & CAISO regulation and reserve price files & Interval-aligned $C_t^{k,s}$ arrays & Drives AS capacity revenue.\\
MPC implementation & Forecasts, prices, BESS state, EV state, tariff parameters & DA/RT decisions, dispatch logs, per-day solver outputs & Source of monthly and daily result tables.\\
Financial post-processing & Dispatch logs and price/tariff data & Monthly summary, daily detail, product-level WM tables & Source of results section and appendix accounting tables.\\
\bottomrule
\end{tabularx}
\end{table*}

A useful way to audit the data pipeline is to ask whether every result in the paper can be traced backward to one of these modules. The monthly financial summary traces to financial post-processing; the representative-day 6-panel plot traces to the per-day solver-output folder; EV service statements trace to the processed EV session table and execution logs; and AS revenue traces to the AS price arrays and awarded capacities. This traceability is especially important because the paper compares forecast and controller choices. Without clear provenance, a difference between controllers could be caused by a changed input file rather than by the controller logic itself.

\subsection{Figures and site documentation}
The manuscript uses two kinds of figures. First, system and workflow figures explain the control architecture, data streams, and market timing. These figures are schematic and should use consistent fonts, colors, and line weights. Second, result figures are generated directly from simulation CSVs and solver-output folders. Monthly figures summarize the controller--forecast--market comparison, daily figures show variability across July 2025, and representative-day figures expose the detailed DA/RT operation.

Additional site-context figures can be included when available: a UCSD campus map identifying the studied charging region, EV charging-station photographs or station-layout maps, and a BESS photograph or site diagram. These are useful for case-study clarity but should not replace quantitative results. In the Overleaf project, the image filenames used by the manuscript should match the uploaded figure filenames exactly.
\section{Financial Accounting and Evaluation Metrics}
\label{sec:metrics}

The financial evaluation is performed on executed decisions. The primary metric is monthly net value:
\begin{align}
V^{\mathrm{month}}=R^{\EV}+R^{\WM}-C^{\TOU}-C^{\PD}-C^{\NCD}.
\label{eq:month_value_metric}
\end{align}
This metric is reported for retail-only and wholesale-enabled modes. Retail-only mode sets wholesale participation to zero and evaluates the EV-BESS controller under retail tariff and EV user-payment terms. Wholesale-enabled mode allows DA/RT energy and AS capacity decisions. Comparing these modes reveals whether wholesale participation improves net value after retail impacts.

Daily metrics are also computed, but they must be interpreted carefully because demand charges are monthly or billing-period quantities. A daily table can assign incremental demand-charge impacts to the day on which a new peak occurs, but the economic meaning is different from an energy charge that accrues independently each interval. For this reason, the paper reports both monthly summaries and daily traces. Monthly summaries provide the final objective comparison; daily traces diagnose when high-impact events occur.

Wholesale revenue is decomposed by product and by asset where possible. The product decomposition separates DA energy, RT energy, regulation up, regulation down, spinning reserve, and non-spinning reserve. The asset decomposition separates BESS and EV contributions using implementation variables and ex-post attribution. This attribution should be interpreted as accounting, not as a separate physical market settlement.

\section{Results}
\label{sec:results}

\subsection{Completed scenario set}
The results are organized around the full scenario matrix:
shrinking versus rolling MPC, perfect versus persistence EV information, and retail-only versus retail+wholesale operation. This organization is important because no single controller--forecast pair should be treated as the only result. Perfect EV information provides an information benchmark, persistence provides an implementable EV forecast, shrinking MPC provides a strong workplace-charging baseline, and rolling MPC provides the adaptable execution architecture.

The numerical financial tables in the submission version should be generated directly from the current monthly and daily financial-summary CSVs. In this draft, the table structure is written so that updated values can be inserted without changing the text. The expected columns are net value, EV user revenue, wholesale revenue, TOU cost, peak demand charge, non-coincident demand charge, and delivered EV energy. All entries should be computed from executed decisions rather than from future decisions.

\subsection{Pseudo-algorithm for the execution loop}
The execution loop is summarized in Algorithm~\ref{alg:mpc_execution}. The same algorithmic skeleton applies to both shrinking and rolling MPC; the horizon-update rule is controller-specific.

\begin{algorithm}[!t]
\caption{Execution-aware EV--BESS MPC simulation}
\label{alg:mpc_execution}
\begin{algorithmic}[1]
\STATE Select controller type and EV-information mode.
\STATE Initialize BESS SOC, EV delivered-energy states, retail demand-charge states, and DA/RT references.
\FOR{each control interval $t$ in the study month}
    \STATE Update realized EV arrivals, EV departures, delivered energy, BESS SOC, and active DA commitments.
    \STATE Build EV information set $\widehat{\Xi}_{t:\calH_t}$ using perfect EV information or persistence EV forecasting.
    \STATE Read known tariff, DA/RT price, and AS price inputs over $\calH_t$.
    \STATE Solve the EV--BESS retail--wholesale MILP with the controller-specific horizon $\calH_t$.
    \STATE Execute only the first 15-minute EV charging and BESS dispatch decisions.
    \STATE Update physical states and financial accumulators from executed quantities.
\ENDFOR
\STATE Aggregate monthly, daily, and representative-day financial and operational metrics.
\end{algorithmic}
\end{algorithm}

The final aggregation step is essential. MPC produces a sequence of actions at each solve, but most future actions are overwritten by later solves. Financial accounting should therefore use executed decisions. This is why the paper uses the phrase execution-aware evaluation.

\IfFileExists{tables/monthly_financial_results.tex}{\input{tables/monthly_financial_results.tex}}{
\begin{table*}[!t]
\centering
\caption{Scenario matrix and financial-accounting columns to be populated from the monthly financial summary CSVs.}
\label{tab:monthly_results}
\scriptsize
\begin{tabularx}{\textwidth}{@{}lllX@{}}
\toprule
Controller & EV information & Market mode & Financial columns to report from CSV \\
\midrule
Shrinking & Perfect & Retail only / retail+wholesale & Net value, EV revenue, WM revenue, TOU cost, peak demand charge, NCD charge, delivered EV energy.\\
Shrinking & Persistence & Retail only / retail+wholesale & Same columns, using persistence EV forecasts.\\
Rolling & Perfect & Retail only / retail+wholesale & Same columns, using the completed rolling-perfect simulations.\\
Rolling & Persistence & Retail only / retail+wholesale & Same columns, using persistence EV forecasts.\\
\bottomrule
\end{tabularx}
\end{table*}
}

\subsection{Monthly value-stack decomposition}
The value-stack decomposition should report at least five components: EV user payments, wholesale revenue, TOU energy cost, peak-demand cost, and non-coincident demand cost. A decomposition is more informative than a single net-value number because it reveals why a case wins or loses. In the completed persistence results, rolling MPC increases wholesale revenue relative to shrinking MPC, but the incremental wholesale gain is smaller than the increase in retail exposure. This is exactly the type of conclusion that would be hidden if only gross wholesale revenue were reported.

\begin{figure*}[!t]
\centering
\MaybeFigure{fig_monthly_value_stack_decomposition.png}{width=0.88\textwidth}
\caption{Monthly value-stack decomposition. Positive components represent EV user revenue and wholesale-market revenue; negative components represent TOU, peak-demand, and non-coincident demand charges.}
\label{fig:value_stack}
\end{figure*}

\begin{figure*}[!t]
\centering
\MaybeFigure{fig_monthly_financial_summary_perfect.png}{width=0.86\textwidth}
\caption{Monthly financial summary for the perfect-EV-information case, comparing retail-only operation with retail+wholesale value stacking. The grouped bars separate wholesale-market revenue, EV user payments, TOU energy cost, peak-demand charge, non-coincident demand charge, and total net value.}
\label{fig:monthly_financial_perfect}
\end{figure*}

\subsection{Daily variability and month-level interpretation}
Monthly values can hide operational variability. A few days can dominate the value of a monthly demand charge, and a few high-price market intervals can dominate AS or energy revenue. The daily financial detail files should therefore be used to show how net value changes across July. This daily view also helps diagnose whether a controller performs poorly because of many small errors or because of a small number of high-impact intervals. Figures~\ref{fig:daily_component_perfect}--\ref{fig:daily_wm_breakdown_perfect} add component-level views for the perfect-EV-information case: the first separates daily retail and EV-payment components, the second displays daily stacked value, and the third isolates wholesale-market energy, capacity, BESS, and EV contributions.

\begin{figure*}[!t]
\centering
\MaybeFigure{fig_daily_total_revenue_comparison.png}{width=0.88\textwidth}
\caption{Daily net-value comparison across controller--forecast scenarios. The figure highlights whether monthly differences arise from systematic daily behavior or from a small number of high-impact operating days.}
\label{fig:daily_revenue}
\end{figure*}

\begin{figure*}[!t]
\centering
\MaybeFigure{fig_daily_component_comparison_perfect.png}{width=0.98\textwidth}
\caption{Daily component-level comparison for the perfect-EV-information case. The five panels separate wholesale-market revenue, TOU energy cost, peak-demand charge, non-coincident demand charge, and EV user-payment revenue, allowing the retail-only and retail+wholesale cases to be compared day by day.}
\label{fig:daily_component_perfect}
\end{figure*}

\begin{figure*}[!t]
\centering
\MaybeFigure{fig_daily_value_stack_perfect.png}{width=0.92\textwidth}
\caption{Daily value-stack decomposition for the perfect-EV-information case. The top panel shows retail-only operation, while the bottom panel shows retail+wholesale operation. Stacked bars report revenue and cost components, and the black line reports daily net value.}
\label{fig:daily_value_stack_perfect}
\end{figure*}

\begin{figure*}[!t]
\centering
\MaybeFigure{fig_daily_wm_breakdown_perfect.png}{width=0.94\textwidth}
\caption{Daily wholesale-market revenue breakdown for the perfect-EV-information retail+wholesale case. The upper panel separates energy and capacity revenue; the lower panel attributes wholesale-market revenue between BESS and EV flexibility.}
\label{fig:daily_wm_breakdown_perfect}
\end{figure*}

\subsection{Representative-day six-panel operational analysis}
A representative-day plot is necessary because the optimization model is multi-layered. Monthly tables show value, but they do not show how EV charging, BESS SOC, wholesale energy, AS capacity, and retail import interact. The six-panel plot from the solver-output folders is used to explain one operating day in detail. Its caption states the date, controller, forecast method, and operating mode. Each panel should then be connected to the model: EV charging shows service and flexibility, BESS power/SOC shows energy shifting and capacity headroom, DA/RT schedules show wholesale participation, AS awards show reserve capacity, and grid import shows retail-tariff exposure.

\begin{figure*}[!t]
\centering
\MaybeFigure{fig_representative_day_6panel.png}{width=0.92\textwidth}
\caption{Representative-day six-panel operation plot selected from the RT solver-output folder. The panels should be read as a detailed operational trace linking EV dispatch, BESS operation, grid import, wholesale energy quantities, AS awards, and financial impacts for one high-information operating day.}
\label{fig:representative_day}
\end{figure*}

A good representative-day discussion should not simply say that the plot is complex. It should explain the causal chain. For example, if wholesale prices or AS prices are high in a particular period, the controller may reserve BESS discharge headroom or reduce EV charging relative to baseline. If EV demand is then shifted into a later period, the retail meter may experience higher import during a TOU or demand-charge window. The six-panel figure should therefore be used to show why the net result can differ from the wholesale result.

\subsection{How the representative day should be selected}
The representative-day figure should not be chosen only because it is visually attractive. It should be chosen because it helps explain a mechanism in the monthly results. A good candidate is a day where wholesale revenue is large but net value is not correspondingly high, or a day where a demand-charge threshold changes the monthly economics. The selection procedure should therefore inspect daily net value, daily wholesale revenue, daily TOU cost, and daily demand-charge increments. The chosen day should be reported in the caption and in the text.

Once the day is selected, the six-panel analysis should answer five questions. First, when do EV sessions create high aggregate available charging capability? Second, when does the BESS charge or discharge, and how does SOC move? Third, when are DA/RT energy positions nonzero? Fourth, when are AS capacity awards scheduled, and are they upward or downward products? Fifth, how do these decisions affect grid import at the retail meter? This structure turns the 6-panel figure from a diagnostic plot into a scientific explanation.

\begin{table*}[!t]
\centering
\caption{Checklist for interpreting the representative-day six-panel plot.}
\label{tab:repday_checklist}
\scriptsize
\begin{tabularx}{\textwidth}{@{}p{0.20\textwidth}p{0.30\textwidth}X@{}}
\toprule
Panel or quantity & What to inspect & How to connect it to the model \\
\midrule
EV charging & Timing, magnitude, and deviation from baseline & Shows whether EV flexibility is used and whether service is preserved.\\
EV capability & Number of connected vehicles and charger-limit envelope & Explains why upward/downward EV flexibility changes through the day.\\
BESS power and SOC & Charging, discharging, and SOC headroom & Shows energy shifting and AS duration feasibility.\\
DA/RT energy & Energy positions and RT deviations & Shows how market schedules interact with execution.\\
AS capacity & RU/RD/SP/NSP awards & Shows whether the site is reserving upward or downward capability.\\
Grid import & Meter power and peak intervals & Connects market actions to TOU and demand-charge exposure.\\
\bottomrule
\end{tabularx}
\end{table*}

\subsection{Summary of results and controller interpretation}
The current evidence supports a balanced interpretation. Shrinking MPC should not be treated as a weaker baseline; it is well suited to daytime workplace sessions that close within the day. Rolling MPC is more general and is expected to become more valuable when the EV fleet has overnight sessions or cross-day flexibility. Perfect forecasts quantify information value, while persistence forecasts quantify a practical controller. Future results should complete the rolling-perfect case and extend the month set before making stronger claims about seasonal robustness.

\section{Discussion}
\label{sec:discussion}

The main scientific message is that execution-aware value stacking is a net-value problem, not a revenue-maximization problem. A behind-the-meter EV-BESS site can access wholesale revenue only through physical decisions that also affect the retail meter. When the model increases BESS charging or changes EV charging to support a market product, it may also increase TOU cost or demand charges. Thus, a controller that is more active in the wholesale market can appear better under gross revenue but worse under total value.

The second message is that controller architecture should match the flexibility structure of the EV data. For the studied workplace data, zero overnight sessions mean that the daily boundary is meaningful. Shrinking MPC naturally respects this boundary. Rolling MPC is not expected to outperform merely because it has a moving horizon; its advantage is adaptability. If future cases include overnight charging, fleet vehicles, residential charging, or stronger cross-day uncertainty, rolling MPC may become more valuable. This distinction should be made clear in the final paper so that the results are not misread as a failure of rolling control.

The third message is methodological. EV flexibility without V2G must be explained carefully. The EV fleet does not sell energy from vehicle batteries. It provides flexibility by changing the timing of charging demand relative to a baseline and by coordinating that timing with BESS operation. The BESS can perform direct charge/discharge energy shifting; the EV fleet provides controllable demand. The combination can be valuable, but only if user energy requirements and all market feasibility constraints are satisfied.

\section{Limitations and Future Work}
\label{sec:limitations}

The present paper reports a July 2025 sample-month case study. The results should be expanded for the submission package or framed explicitly as a sample-month demonstration. Future work should simulate additional months and seasons, evaluate tariff sensitivity, complete rolling-perfect results, test different RT time limits and optimality gaps, revise and validate the EV baseline calculation, integrate PV and building load into the active model, and evaluate the advanced EV forecast module. Overnight charging scenarios are particularly important because they test the hypothesis that rolling MPC becomes more valuable when EV flexibility crosses daily boundaries.

The current model also uses expected deployment factors for AS energy accounting. A future version can compare expected-deployment accounting with realized deployment traces or scenario-based deployment. Similarly, the current model treats the site as a market participant with the ability to map behind-the-meter flexibility into wholesale energy and AS products. A final deployment study should include market-registration details, telemetry requirements, minimum bid sizes, aggregation rules, and settlement assumptions.

\section{Conclusion}
\label{sec:conclusion}

This paper develops a detailed EV-BESS value-stacking framework for a campus behind-the-meter system under retail--wholesale coupling. The formulation co-optimizes EV charging, BESS operation, retail energy and demand charges, EV user payments, DA/RT wholesale energy, and AS capacity products. The paper compares shrinking and rolling MPC and distinguishes perfect and persistence EV forecasts. The July 2025 results show that shrinking and rolling controllers should be interpreted as complementary frameworks. In the present daytime workplace scenario, shrinking persistence produces higher net value than rolling persistence in the wholesale-enabled case even though rolling persistence earns more wholesale revenue. This outcome highlights the central lesson of the study: more market activity is not automatically better when retail tariff exposure and EV service obligations are included. The framework is designed to be extended to additional months, PV and building load integration, additional EV-forecasting methods and overnight charging scenarios.

\appendix
\section{Complete MILP Formulation}
\label{app:complete_milp}

This appendix records the implementation-level constraints. They are separated from the main text to keep the main formulation readable while preserving reproducibility.

\subsection{EV availability and energy state}
\begin{align}
b_{t,i}^{\EV} &= \begin{cases}1, & a_i\le t\le d_i,\\0, & \text{otherwise},\end{cases}\label{eq:app_ev_avail}\\
e_{a_i,i} &= e_{a_i},\qquad e_{0,i}=0,\label{eq:app_ev_init}\\
e_{t,i} &= e_{t-1,i}+p_{t,i}^{\EV}\DeltaT/\eta_i^{\EV},\label{eq:app_ev_dyn}\\
0&\le e_{t,i}\le e_i^{\req},\qquad
0\le p_{t,i}^{\EV}\le \hat P^{\EV}b_{t,i}^{\EV},\label{eq:app_ev_bounds}\\
e_{d_i,i}&\ge e_i^{\req}.\label{eq:app_ev_req}
\end{align}
For rolling MPC, service closure constraints are applied so that energy required for currently relevant same-day sessions cannot be pushed beyond the valid service window. The precise mask depends on whether the horizon is a same-day shrinking horizon or a rolling cross-day horizon.

\subsection{BESS SOC, split, and terminal equality}
\begin{align}
p_t^{\ch,\BESS}&=p_t^{\ch,\BESS,\WM}+p_t^{\ch,\BESS,\NWM},\label{eq:app_bess_chsplit}\\
p_t^{\dch,\BESS}&=p_t^{\dch,\BESS,\WM}+p_t^{\dch,\BESS,\NWM},\label{eq:app_bess_dchsplit}\\
p_t^{\BESS}&=p_t^{\ch,\BESS}-p_t^{\dch,\BESS},\label{eq:app_bess_power}\\
\soc_t^{\BESS}&=\soc_{t-1}^{\BESS}+\frac{\DeltaT}{C^{\BESS}}
\left(p_t^{\ch,\BESS}\sqrt{\gamma}-\frac{p_t^{\dch,\BESS}}{\sqrt{\gamma}}\right),\label{eq:app_bess_soc}\\
\soc_{\min}^{\BESS}&\le \soc_t^{\BESS}\le \soc_{\max}^{\BESS},\label{eq:app_bess_bounds}\\
\soc_{T_d}^{\BESS}&=\soc_{0,d}^{\BESS}.\label{eq:app_bess_terminal}
\end{align}
The wholesale/non-wholesale split is used for accounting and attribution. The physical SOC equation uses the total charge and discharge powers.

\subsection{Direction binaries}
\begin{align}
0&\le p_t^{\ch,\BESS}\le P_{\max}^{\BESS}u_t^{\ch,\BESS},\label{eq:app_bess_chbin}\\
0&\le p_t^{\dch,\BESS}\le P_{\max}^{\BESS}u_t^{\dch,\BESS},\label{eq:app_bess_dchbin}\\
u_t^{\ch,\BESS}+u_t^{\dch,\BESS}&\le 1.\label{eq:app_bess_dirbin}
\end{align}
Analogous direction binaries are used for EV net load-increase/load-reduction variables and for net wholesale charge/discharge accounting. These binaries prevent the model from claiming simultaneous upward and downward capability from the same physical margin.

\subsection{EV net-deviation and big-M formulation}
The EV fleet is decomposed relative to the baseline:
\begin{align}
p_t^{\EV}-B_t^{\EV}&=p_t^{\ch,\EV}-p_t^{\dch,\EV},\label{eq:app_ev_deviation}\\
0&\le p_t^{\ch,\EV}\le M^{\EV}u_t^{\ch,\EV},\label{eq:app_ev_ch}\\
0&\le p_t^{\dch,\EV}\le M^{\EV}u_t^{\dch,\EV},\label{eq:app_ev_dch}\\
u_t^{\ch,\EV}+u_t^{\dch,\EV}&\le 1,\label{eq:app_ev_dir}\\
M^{\EV}&=2P_{t,\max}^{\EV}.\label{eq:app_ev_bigM}
\end{align}
Here $p_t^{\dch,\EV}$ is a load-reduction variable, not physical EV discharge. It measures how much charging is reduced relative to the baseline.

\subsection{Wholesale capacity envelopes}
The total upward and downward deliverable capabilities are
\begin{align}
\hat p_t^{\up}&=P_{\max}^{\BESS}+B_t^{\EV},\label{eq:app_upcap}\\
\hat p_t^{\down}&=P_{\max}^{\BESS}+P_{t,\max}^{\EV}-B_t^{\EV}.\label{eq:app_downcap}
\end{align}
The DA/RT energy positions and AS awards must satisfy
\begin{align}
&(p_t^{\DA})^+ +(p_t^{\RT})^+
+c_t^{\RU,\DA}+c_t^{\RU,\RT}
+c_t^{\SP,\DA}+c_t^{\SP,\RT}\nonumber\\
&\qquad +c_t^{\NSP,\DA}+c_t^{\NSP,\RT}
\le \hat p_t^{\up},\label{eq:app_upenv}\\
&(p_t^{\DA})^- +(p_t^{\RT})^-
+c_t^{\RD,\DA}+c_t^{\RD,\RT}
\le \hat p_t^{\down}.\label{eq:app_downenv}
\end{align}
Positive and negative components are represented with auxiliary variables in the implementation rather than nonlinear max functions.

\subsection{Net power flow and AS deployment accounting}
Expected deployment creates additional net charge/discharge requirements:
\begin{align}
q_t^{\deploy}=&\alpha_t^{\RU}(c_t^{\RU,\DA}+c_t^{\RU,\RT})-
\alpha_t^{\RD}(c_t^{\RD,\DA}+c_t^{\RD,\RT})\nonumber\\
&+\alpha_t^{\SP}(c_t^{\SP,\DA}+c_t^{\SP,\RT})+
\alpha_t^{\NSP}(c_t^{\NSP,\DA}+c_t^{\NSP,\RT}).\label{eq:app_deploy}
\end{align}
The net wholesale movement combines DA/RT energy and expected AS deployment. Direction binaries and auxiliary variables convert this signed movement into nonnegative charge and discharge components. This is the implementation-level reason why AS awards can affect both wholesale revenue and BESS/EV feasibility.

\subsection{AS duration constraints}
Capacity awards must be deliverable for the required duration $T_{\mathrm{DU}}$:
\begin{align}
\soc_t^{\BESS}-\soc_{\min}^{\BESS}
&\ge \frac{T_{\mathrm{DU}}}{C^{\BESS}\sqrt{\gamma}}
\left(c_t^{\RU,\tot}+c_t^{\SP,\tot}+c_t^{\NSP,\tot}\right),\label{eq:app_duration_up}\\
\soc_{\max}^{\BESS}-\soc_t^{\BESS}
&\ge \frac{T_{\mathrm{DU}}\sqrt{\gamma}}{C^{\BESS}}c_t^{\RD,\tot},\label{eq:app_duration_down}\\
c_t^{k,\tot}&=c_t^{k,\DA}+c_t^{k,\RT}.\label{eq:app_capacity_total}
\end{align}
The present implementation uses $T_{\mathrm{DU}}=0.5$ h. EV energy-duration limits can be added analogously when product-level EV duration accounting is activated.

\subsection{RT reference tracking and penalties}
In RT, DA reference values are not ignored. The implemented model penalizes deviations from DA reference schedules:
\begin{align}
\rho_t^{\WM}=&M_t^{\WM}\left(|p_t^{\DA}-\bar p_t^{\DA}|+
\sum_{k\in\calK}|c_t^{k,\DA}-\bar c_t^{k,\DA}|\right),\label{eq:app_rt_penalty}\\
M_t^{\WM}=&M_0|C_t^{\mathrm{LMP},\RT}|.\label{eq:app_rt_weight}
\end{align}
This soft-tracking representation matches the current simulation implementation better than a hard-fix statement. It keeps RT flexible while discouraging unrealistic DA/RT inconsistency.

\subsection{Revenue decomposition}
Wholesale revenue is decomposed as
\begin{align}
R^{\WM,\energy}&=\sum_t\DeltaT\left(C_t^{\mathrm{LMP},\DA}p_t^{\DA}+C_t^{\mathrm{LMP},\RT}p_t^{\RT}\right)+R^{\mathrm{deploy}},\label{eq:app_energy_rev}\\
R^{\WM,\capacity}&=\sum_t\DeltaT\sum_{k\in\calK}\sum_{s\in\calS}C_t^{k,s}c_t^{k,s},\label{eq:app_capacity_rev}\\
R^{\WM}&=R^{\WM,\energy}+R^{\WM,\capacity}.\label{eq:app_total_rev}
\end{align}
Asset-level EV/BESS attribution is an ex-post accounting calculation based on EV and BESS wholesale variables. It should not be interpreted as a separate physical constraint.




\balance
\bibliographystyle{elsarticle-num}
\bibliography{references}

\end{document}
